\documentclass{acm_proc_article-sp}

\usepackage{graphicx,url}
\usepackage[portuges]{babel}   
\usepackage[utf8]{inputenc}
\usepackage{paralist}


\newcommand{\fig}[4][htb]{
  \begin{figure}[#1]
    {\centering{\includegraphics[#4]{fig/#2}}\par}
%    \caption{#3}
    \label{fig:#2}
  \end{figure}
}

\newcommand{\tab}[4][h]{
  \begin{table}[h]
    {\centering\footnotesize\textsf{\input{fig/#2.tab}}\par}
%    \caption{#3}
    \label{tab:#2}
  \end{table}
}

\newcommand{\emote}{\textsc{EPOSMote}}
\newcommand{\epos}{\textsc{Epos}}
\newcommand{\eplc}{\textsc{EposMotePLC}}

\begin{document}
	
	\title{
	  {\centering\textwidth}
	  \begin{tabular}[t]{ccc}
	    \raisebox{-0.7cm}[0pt][0pt]{\resizebox*{!}{1.5cm}{\includegraphics{fig/ufsc}}}
	    &
	    \textbf{\Large Universidade Federal de Santa Catarina}
	    &
	    \raisebox{-0.7cm}[0pt][0pt]{\resizebox*{!}{1cm}{\includegraphics{fig/logolisha}}}\\
	    &\textbf{\Large Laboratório de Integração Software/Hardware} &\\
	    &\textbf{\Small http://www.lisha.ufsc.br --- Fone/Fax: +55 48 3721-9516} &
	 \end{tabular}
	 \newline
	 \newline
	 {\LARGE \bf Smart Cities - Viabilizando Cidades Inteligentes}
	}

\maketitle

\section*{}
A rápida evolução da tecnologia tem ocasionado uma mudança expressiva no
ambiente urbano. Há uma crescente instrumentação de áreas urbanas com sensores
interconectados por redes sem fio que fornecem diversos dados sobre o
ambiente. O desenvolvimento de sistemas capazes de tirar proveito das novas
tecnologias de informação e comunicação e que façam uso da massa de dados
disponível, agregando-lhes significado e integrando-as de uma forma sensível
ao contexto em que são produzidos, será capaz de prover novos níveis de
``inteligência urbana'', viabilizando a concepção de Cidades Inteligentes
(smart cities).

\fig{smartgrid2}{}{width=\columnwidth}

Uma cidade inteligente necessita de soluções para vários problemas, como
trânsito, vigilância, atendimento de emergência, monitoramento ambiental,
saúde, educação e inclusão digital, bem como permite uma maior eficiência e
transparência na gestão pública. Para que os objetivos das cidades
inteligentes sejam alcançados é preciso fomentar o conceito de computação
ubíqua, no qual a presença de tecnologia computacional está presente no
dia-a-dia das pessoas, seja nos objetos, ambientes, ferramentas, veículos,
equipamentos e edificações.

\section*{Viabilizando Cidades Inteligentes}

Para viabilizar o paradigma da computação ubíqua, vários desafios precisam ser suplantados. Inicialmente, o ambiente físico precisa ser instrumentado por intermédio do uso de dispositivos sensores. Tais sensores devem ser conectados e trabalharem de forma colaborativa, formando uma rede de sensores. Logo, modelos de armazenamento, tratamento e disseminação de dados específicos para redes de sensores urbanas devem ser investigados de modo a prover soluções customizadas para esses ambientes, respeitando o requisito de eficiência em energia de tais redes.

\fig{cia2-overview}{}{width=\columnwidth}

\section*{Internet das Coisas} 

Um dos pilares de Cidades Inteligentes é a integração de Redes de Sensores sem
Fio com objetos cotidianos, que passariam então a interagir entre si e com a
infraestrutura das Cidades Inteligentes para nos propiciar comodidade e
segurança, traz a tona o conceito de Internet das Coisas (do inglês Internet
of things – IoT). De forma a viabilizar a implementação de comunicação entre
milhares de objetos é fundamental que novos protocolos sejam estudados, uma
vez que as tecnologias atuais demandam um capacidade computacional e
energética incompatíveis com o conceito de IoT.

\fig{iot}{}{width=\columnwidth}

\section*{Smart Grids}

Smart grid é o uso de tecnologia da informação para o provimento de energia em
uma Cidade Inteligente. SmartGrids envolvem a instalação de sensores nas
linhas da rede de energia elétrica estabelecendo uma comunicação confiável em
dua vias com os dispositivos que utilizam a rede. Desta forma, é possivel que
a operação e o desempenho da rede sejam sentidos pela concesionária de
energia, que pode agir para que o sistema disponibilize-a da maneira mais
eficiente. De fato, com o uso desta tecnologia se vislumbra um futuro em que a
própria geração de energia possa ser realizada de forma distribuida, com o uso
de energias limpas, utilizando o sistema de tecnologia da informação para
coordenar de uma forma confiável o sistema.

\section*{Smart Spaces}

Smart Space é a aplicação do conceito de ambientes inteligentes no espaço
público de uma cidade. Aplicações para SmartSpaces envolvem a distribuição de
uma rede de sensores sem fio responsáveis por coletar informações do ambiente
e implementar serviços inteligentes aos cidadões. Exemplos de aplicações em
SmartSpaces envolve aplicações como monitoramento de rios, monitoramento de
trânsito, monitoramento de condições ambientais, segurança pública e o
fornecimentos de serviços tais como alertas e informativos aos cidadões.

O LISHA atualmente esta participando de uma rede de instituições que irá investigar diversas tecnologias para Cidades Inteligentes, realizando experimentações em aplicações de sensoriamento de rios e de tráfego urbano.

\end{document}

